\bookchapter{Capítulo 14: Tema de Gerenciadores de Senhas}

\chapterintro{Na era digital, a senha é a chave da porta da sua casa. Mas e quando as chaves são muitas e complexas demais? O gerenciador de senhas é o seu “chaveiro mestre universal”.}

\booksection{I. Por que Você Precisa de um Gerenciador de Senhas?}

Em média, uma pessoa moderna tem dezenas ou até centenas de contas online. Se cada conta usar a mesma senha, assim que um site for invadido, todas as contas correm perigo. Se cada conta usar uma senha diferente e complexa, é simplesmente impossível decorar todas.

O gerenciador de senhas resolve esse dilema: você só precisa memorizar uma senha mestra, e ele lembra todas as outras para você.

\booksection{II. Como Funciona um Gerenciador de Senhas}

\begin{enumerate}
  \item \textbf{Armazenamento criptografado}: todas as senhas são armazenadas com criptografia de alto nível, e só você pode descriptografá-las.
  \item \textbf{Senha mestra}: você define uma senha mestra para desbloquear o cofre de senhas. A senha mestra em si não é armazenada no servidor.
  \item \textbf{Preenchimento automático}: na tela de login de sites ou apps, o gerenciador de senhas preenche automaticamente usuário e senha.
  \item \textbf{Geração de senhas}: pode gerar senhas aleatórias e complexas, aumentando a segurança.
\end{enumerate}

\booksection{III. Introdução aos Princípios de Criptografia}

Entender os princípios de criptografia do gerenciador de senhas ajuda a avaliar se um provedor é confiável. A seguir, explico alguns conceitos essenciais em linguagem acessível.

\booksubsection{Criptografia AES-256}

O AES (Advanced Encryption Standard, Padrão de Criptografia Avançada) é um algoritmo de criptografia simétrica reconhecido pelo NIST (Instituto Nacional de Padrões e Tecnologia dos EUA), amplamente usado por bancos, forças armadas e governos. O “256” refere-se ao comprimento da chave, de 256 bits — teoricamente, seria necessário mais tempo do que a idade do universo para quebrá-la por força bruta. Os principais gerenciadores de senhas (1Password, Bitwarden, KeePass, etc.) usam AES-256 para criptografar seu cofre de senhas.

\booksubsection{Arquitetura de Conhecimento Zero (Zero-Knowledge Architecture)}

“Conhecimento zero” não significa que o provedor não armazena nada, mas sim que\textbf{a criptografia e a descriptografia são feitas localmente no seu dispositivo; o servidor apenas guarda o texto cifrado já criptografado}. O fluxo específico:

\begin{enumerate}
  \item Você digita a senha mestra no cliente.
  \item O cliente deriva a chave de criptografia a partir da senha mestra (essa etapa é realizada no seu celular/computador).
  \item Apenas o texto cifrado é enviado ao servidor.
  \item No próximo login, o servidor envia o texto cifrado de volta, e o cliente descriptografa novamente com a senha mestra.
\end{enumerate}

Ponto-chave:\textbf{a senha mestra nunca sai do seu dispositivo}; ela também não está no servidor. Portanto, mesmo que o provedor seja invadido, o hacker só obtém um monte de texto cifrado indecifrável — a menos que consiga adivinhar sua senha mestra.

\begin{tcbinfo}{Por que "conhecimento zero" significa que nem o provedor consegue ver suas senhas?}
Porque a chave de descriptografia é derivada da sua senha mestra, e a senha mestra só existe na sua mente e só é usada no seu dispositivo local. No servidor do provedor, não há nem a senha mestra, nem a chave de descriptografia — apenas texto cifrado. É por isso que provedores sérios como Bitwarden e 1Password declaram explicitamente que “nós também não podemos ver suas senhas” — não é que eles não queiram ver, é que tecnicamente é impossível.

\end{tcbinfo}

\booksubsection{Funções de Derivação de Chave: PBKDF2 e Argon2}

A senha mestra não pode ser usada diretamente como chave de criptografia — ela precisa passar por uma “função de derivação de chave”. O papel dessa função é\textbf{tornar o cálculo deliberadamente mais lento}, elevando o custo da quebra por força bruta a um nível inviável.

\begin{itemize}
  \item \textbf{PBKDF2}: padrão veterano, amplamente usado pelo iCloud Keychain, Bitwarden, LastPass e outros. Desacelera a quebra por meio de iterações repetidas de hash (geralmente acima de 100 mil vezes).
  \item \textbf{Argon2}: vencedor do concurso de hash de senhas de 2015, mais moderno que o PBKDF2. Além do "consumo de poder computacional", introduz também o "consumo de memória" — dificultando a aceleração da quebra até mesmo por hardware dedicado como GPU/ASIC. Bitwarden, KeePassXC e outros já oferecem suporte.
\end{itemize}

Quanto maior o número de iterações, mais seguro, mas o desbloqueio também fica mais lento. Provedores sérios divulgam publicamente o algoritmo de derivação e o número de iterações que utilizam, para facilitar a auditoria pela comunidade de segurança.

\booksubsection{Alerta de Caso: O Vazamento de Dados do LastPass em 2022}

Em agosto de 2022, o veterano gerenciador de senhas LastPass anunciou que seu ambiente de desenvolvimento havia sido invadido; em novembro do mesmo ano, divulgou ainda que os hackers haviam levado os backups dos cofres de senhas dos usuários (texto cifrado). As principais lições do incidente:

\begin{itemize}
  \item \textbf{A arquitetura de conhecimento zero segurou a linha de defesa}: como o cofre em si é criptografado, os hackers obtiveram o texto cifrado, mas não conseguiram decifrá-lo — desde que a senha mestra do usuário fosse forte o suficiente.
  \item \textbf{Usuários com senha mestra fraca são os mais vulneráveis}: os hackers podem fazer quebra por força bruta offline das senhas mestras fracas; portanto, usuários que usam senha mestra fraca ou reutilizada correm risco real.
  \item \textbf{Os metadados também vazaram}: embora o conteúdo das senhas esteja criptografado, os metadados, como URLs dos sites e anotações, podem não estar criptografados, expondo quais serviços o usuário utilizava.
\end{itemize}

\begin{tcbwarn}{O que Aprendemos com o Caso LastPass}
Primeiro, a senha mestra deve ser longa e única (recomenda-se 16 caracteres ou mais, incluindo maiúsculas, minúsculas, números e símbolos). Segundo, prefira provedores que suportem Argon2id e com alto número de iterações. Terceiro, para contas de alto valor (banco, e-mail), ative obrigatoriamente uma segunda verificação independente; não dependa apenas do gerenciador de senhas.

\end{tcbwarn}

\booksection{IV. Principais Gerenciadores de Senhas Recomendados}

\booksubsection{Nativos do sistema:}

\begin{itemize}
  \item \textbf{iCloud Keychain}: primeira escolha para usuários do ecossistema Apple. Integração perfeita com Safari e apps, gratuito, seguro. A desvantagem é a dificuldade de uso entre plataformas.
  \item \textbf{Gerenciador de Senhas do Google}: prático para usuários Android. Integrado ao Chrome e ao sistema Android, gratuito.
  \item \textbf{Cofre de Senhas Huawei}: nativo do HarmonyOS, adequado para usuários Huawei.
\end{itemize}

\booksubsection{Ferramentas profissionais de terceiros:}

\begin{itemize}
  \item \textbf{Bitwarden}: open source e gratuito, funcionalidades poderosas, suporta todas as plataformas. Altamente recomendado para usuários que valorizam a privacidade.
  \item \textbf{1Password}: interface elegante, funções ricas, ideal para famílias e equipes. Pago.
  \item \textbf{LastPass}: gerenciador de senhas veterano, versão gratuita com funções limitadas.
  \item \textbf{KeePass}: armazenamento totalmente local, open source e gratuito. Ideal para usuários geek.
\end{itemize}

\booksection{V. Avaliação Real da Força da Senha}

Os gerenciadores de senha costumam exibir uma barra indicadora de “força da senha” (vermelho/amarelo/verde), mas essa avaliação geralmente considera apenas o comprimento e os tipos de caracteres, não sendo abrangente. Para julgar de verdade se uma senha é segura, é preciso considerar a “entropia” e se “já foi vazada”.

\booksubsection{Entropia (Entropy): A Medida Científica da Força da Senha}

A “entropia” mede a incerteza de uma senha, em bits. Quanto maior a entropia, maior o número de tentativas necessárias para quebrá-la por força bruta.

\begin{itemize}
  \item \textbf{"123456"}: entropia extremamente baixa, quebrada instantaneamente.
  \item \textbf{"P@ssw0rd"}: parece complexa, mas está em todas as listas de ataque por dicionário; a entropia continua baixa.
  \item \textbf{"correct horse battery staple"}: quatro palavras aleatórias, entropia de cerca de 44 bits, muito mais difícil de quebrar do que "P@ssw0rd".
  \item \textbf{16 caracteres aleatórios (incluindo maiúsculas/minúsculas + números + símbolos)}: entropia de cerca de 100 bits, impossível de quebrar por força bruta com o poder computacional atual.
\end{itemize}

Conclusão:\textbf{“comprimento + aleatoriedade” é mais importante do que “símbolos complexos”}. É por isso que as senhas longas e aleatórias geradas pelo gerenciador de senhas são realmente seguras.

\booksubsection{HaveIBeenPwned: Verifique se Sua Senha Já Foi Vazada}

Mantido pelo pesquisador de segurança Troy Hunt, o\href{https://haveibeenpwned.com}{haveibeenpwned.com}reúne contas e senhas que vazaram em sucessivos grandes incidentes de vazamento de dados. Você pode:

\begin{enumerate}
  \item Digitar seu e-mail e verificar se sua conta foi exposta em algum vazamento.
  \item Usar a função Pwned Passwords para consultar se uma determinada senha está na base de dados de senhas vazadas conhecidas (os gerenciadores de senha e navegadores já têm essa consulta integrada).
\end{enumerate}

Se uma senha aparecer na base de dados de vazamentos,\textbf{não importa o quão complexa ela pareça, deve ser trocada imediatamente}— porque o ataque por dicionário tentará essas senhas conhecidas primeiro.

\begin{tcbinfo}{Princípio da Consulta k-Anonymity}
Preocupado em digitar sua senha no HaveIBeenPwned e ela vazar de novo? Não precisa. As consultas modernas usam o protocolo “k-anonimato”: o navegador envia apenas os 5 primeiros caracteres do hash da senha para o servidor; o servidor retorna todos os hashes que começam com esses 5 caracteres; a comparação é feita localmente. Ou seja, o servidor nunca vê o hash completo da sua senha. É assim que funciona a “detecção de senhas vazadas” do 1Password e do Bitwarden.

\end{tcbinfo}

\booksubsection{Funções de Monitoramento Integradas dos Gerenciadores de Senhas}

\begin{itemize}
  \item \textbf{1Password Watchtower}: monitora em tempo real senhas fracas, senhas repetidas, senhas vazadas, senhas expiradas e sites sem 2FA, exibindo com classificação por cores vermelho/laranja/amarelo.
  \item \textbf{Bitwarden Vault Health}: a versão gratuita já permite ver senhas fracas e repetidas; a versão paga detecta senhas vazadas (com base no HaveIBeenPwned).
  \item \textbf{"Sugestões de Segurança" do iCloud Keychain}: Ajustes do iOS Senhas Sugestões de Segurança, indica quais senhas aparecem em vazamentos conhecidos e quais estão sendo reutilizadas.
  \item \textbf{"Verificador de Senhas" do Gerenciador de Senhas do Google}: integrado ao Chrome, verifica em um toque senhas vazadas, fracas e repetidas.
\end{itemize}

\begin{tcbpractice}{Verifique Agora: Suas Senhas Estão Seguras?}
\begin{enumerate}
  \item Abra o haveibeenpwned.com, digite seu e-mail mais usado e veja em quais incidentes de vazamento ele aparece.
  \item Usuários de iPhone: Ajustes Senhas Recomendações de segurança, trate os avisos vermelhos um por um.
  \item Usuários de Chrome: Configurações Senhas Verificar senhas, modifique conforme as sugestões.
  \item Usuários de 1Password/Bitwarden: Abra o Watchtower/Health Reports, priorize os itens "Vazados" e "Duplicados".
\end{enumerate}

\end{tcbpractice}

\booksection{VI. Passkey: O futuro do login sem senha}

De 2024 a 2026, a tecnologia de senhas está passando por uma revolução — o\textbf{Passkey (Chave de Acesso)}promete eliminar completamente as senhas tradicionais. Apple, Google e Microsoft estão impulsionando juntos, com o endosso da FIDO Alliance e do W3C (World Wide Web Consortium). Esta é a maior mudança na forma de login em décadas.

\booksubsection{O que é Passkey?}

O Passkey é baseado no padrão FIDO2 / WebAuthn, substituindo senhas por um\textbf{par de chaves no dispositivo}. O princípio pode ser simplificado assim:

\begin{itemize}
  \item Seu telefone gera um par de chaves:\textbf{chave pública}(registrada no site) +\textbf{chave privada}(armazenada permanentemente no chip de segurança do seu dispositivo).
  \item Ao fazer login, o site envia um "desafio"; seu telefone assina com a chave privada, e o site verifica com a chave pública.
  \item Antes de assinar, você precisa confirmar sua identidade por biometria (Face ID / impressão digital) ou PIN do dispositivo.
  \item Durante todo o processo,\textbf{nenhuma senha é transmitida pela rede}.
\end{itemize}

\booksubsection{Passkey vs Senhas tradicionais}

\begin{center}
\begin{tabularx}{\linewidth}{|X|X|X|}
\hline
\textbf{Dimensão} & \textbf{Senha tradicional} & \textbf{Passkey} \tabularnewline
\hline
Anti-phishing & Ruim, sites de phishing podem roubar a senha & Naturalmente anti-phishing, a chave privada só assina para o domínio real \tabularnewline
Impacto de vazamento no servidor & O hash da senha pode ser quebrado & O servidor só armazena a chave pública, vazamento sem impacto \tabularnewline
Conveniência & Precisa memorizar/digitar/gerenciar & Login em um segundo com Face ID/impressão digital \tabularnewline
Ataque de repetição & Pode ser interceptado e repetido por atacante intermediário & Cada assinatura é única, não pode ser repetida \tabularnewline
\hline
\end{tabularx}
\end{center}

\booksubsection{Suporte nas plataformas}

\begin{itemize}
  \item \textbf{Apple iCloud Keychain Passkey}: Suporte nativo no iOS 16+ / macOS Ventura+. As Passkeys podem ser sincronizadas entre todos os dispositivos com o mesmo Apple ID. Ao fazer login em um novo dispositivo, não é necessário recriar.
  \item \textbf{Google Password Manager Passkey}: Suporte nativo no Android + Chrome. As Passkeys ficam armazenadas na conta Google, com sincronização entre dispositivos Android. Desde 2023, as contas pessoais do Google já suportam login com Passkey.
  \item \textbf{Microsoft}: Windows 11 / Edge já suportam Passkey. Conta Microsoft, Outlook, Office e outros serviços estão liberando gradualmente o login por Passkey, com gerenciamento pelo Microsoft Authenticator.
  \item \textbf{1Password / Bitwarden}: Já suportam criação, armazenamento e sincronização multiplataforma de Passkeys, libertando você de um único ecossistema.
\end{itemize}

\begin{tcbinfo}{A sincronização entre dispositivos é o avanço chave}
O maior ponto fraco inicial das Passkeys era a “vinculação a um único dispositivo” — trocar de telefone significava recriar tudo. A partir de 2023, Apple, Google e Microsoft implementaram a sincronização em nuvem das Passkeys: iCloud Keychain sincroniza entre dispositivos Apple, Google Password Manager sincroniza entre dispositivos Android. Soluções de terceiros como 1Password permitem a sincronização entre ecossistemas. Isso significa que as Passkeys já têm a praticidade necessária para substituir as senhas.

\end{tcbinfo}

\booksubsection{Como criar e usar Passkey no celular}

\textbf{Caminho no iOS (iPhone):}

\begin{enumerate}
  \item Na página de segurança da conta de um site compatível com Passkey (como Google, Microsoft, Amazon, eBay, PayPal), selecione "Adicionar chave de acesso / Add passkey".
  \item O sistema exibirá a janela de verificação Face ID / Touch ID. Após a confirmação, a Passkey é gerada automaticamente e salva no iCloud Keychain.
  \item No próximo login nesse site, digite o nome de usuário, selecione "Fazer login com iPhone" e passe o Face ID.
  \item Entrada de gerenciamento: Ajustes Senhas encontre o site correspondente para ver a entrada Passkey.
\end{enumerate}

\textbf{Caminho no Android:}

\begin{enumerate}
  \item Na página de segurança da conta de um site compatível com Passkey, selecione "Adicionar chave de acesso".
  \item O Google Password Manager exibirá um pop-up selecione sua conta Google verificação de bloqueio de tela (impressão digital/padrão/PIN).
  \item A Passkey é salva automaticamente no Google Password Manager, com sincronização entre dispositivos Android da mesma conta.
  \item No próximo login, selecione "Usar Passkey" e desbloqueie a tela.
  \item Entrada de gerenciamento: Ajustes do Android Senhas, chaves e preenchimento automático Gerenciador de senhas google.com.
\end{enumerate}

\booksubsection{Login com Passkey entre dispositivos}

Se você estiver fazendo login em um site no Windows, mas a Passkey estiver apenas no iPhone, pode autorizar por QR code:

\begin{enumerate}
  \item Na página de login do site, selecione "Usar Passkey".
  \item Selecione "Usar outro dispositivo".
  \item O computador exibirá um QR code; escaneie com o iPhone.
  \item O iPhone mostrará a verificação Face ID; após confirmar, o login no computador será concluído.
\end{enumerate}

Esse processo utiliza comunicação criptografada dupla via Bluetooth + Internet. A chave privada da Passkey nunca sai do seu iPhone.

\begin{tcbpractice}{Experimente agora: login sem senha}
\begin{enumerate}
  \item Faça login na conta Google Segurança Chaves de acesso Adicionar chave de acesso, experimente o login com Face ID/impressão digital.
  \item Faça login na conta Microsoft Segurança Opções avançadas de segurança Adicionar chave de acesso.
  \item Experimente em serviços que já suportam Passkey, como Taobao, JD.com, eBay, PayPal.
  \item Mantenha a senha tradicional como backup, migrando gradualmente as contas de alto valor para Passkey.
\end{enumerate}

\end{tcbpractice}

\begin{tcbwarn}{Sugestões para o período de transição}
O Passkey ainda não substituiu completamente as senhas — muitos sites ainda estão na fase de coexistência “Passkey + senha”. Sugestão: ative o Passkey prioritariamente para contas de alto valor, mas\textbf{não exclua imediatamente a senha tradicional}, mantenha como backup. Continue usando o gerenciador de senhas, pois nem todos os sites suportam Passkey. Veja o Passkey como um “complemento poderoso ao gerenciador de senhas”, não como uma “substituição imediata”.

\end{tcbwarn}

\booksection{VII. Chave de segurança por hardware}

Se suas exigências de segurança estão acima do usuário comum — como ter muitas criptomoedas, ser executivo/jornalista/ativista, ou proteger contas corporativas essenciais — a chave de segurança por hardware é uma opção mais confiável que o código SMS.

\booksubsection{O que é uma chave de segurança por hardware?}

Uma chave de hardware é um dispositivo físico (semelhante a um pen drive), com um chip de segurança interno que armazena sua chave privada. Requer contato físico ou NFC para concluir a verificação. Produtos representativos:

\begin{itemize}
  \item \textbf{YubiKey}(empresa Yubico): A mais popular, suporta FIDO2/WebAuthn, U2F, OTP, PIV, OpenPGP, com interfaces USB-C / USB-A / NFC. Série 5: 50-70 dólares.
  \item \textbf{Google Titan}: Produzido pelo Google, USB-C + NFC / Bluetooth, cerca de 30-40 dólares.
  \item \textbf{OnlyKey}: Alternativa de código aberto, suporta vários protocolos.
  \item \textbf{Alternativas nacionais}: Como Feitian ePasskey, com preços mais acessíveis.
\end{itemize}

\booksubsection{Por que a chave de hardware é mais segura que o código SMS?}

\begin{itemize}
  \item \textbf{Anti-phishing}: Assim como o Passkey, a chave de hardware só responde ao domínio real. Sites de phishing não conseguem enganá-la para obter o código.
  \item \textbf{Anti-SIM swap}: Códigos SMS podem ser interceptados por ataques de SIM swap. A chave de hardware não precisa de número de telefone.
  \item \textbf{Anti-ataque intermediário (MITM)}: O protocolo FIDO2 inclui um valor de desafio em cada assinatura, impossível de repetir.
  \item \textbf{Isolamento físico}: A chave privada permanece sempre dentro do chip de hardware. Mesmo que um malware comprometa seu computador, não consegue lê-la.
\end{itemize}

\booksubsection{Cenários de uso}

\begin{itemize}
  \item Contas de exchanges de criptomoedas, recuperação de carteira.
  \item SSO corporativo, VPN de trabalho remoto, login SSH em servidores.
  \item "Contas principais" como e-mail, GitHub, redes sociais — se comprometidas, podem redefinir todas as outras contas.
  \item Pessoas de alto risco como jornalistas, advogados, executivos.
\end{itemize}

\booksubsection{Como configurar a chave de hardware como segundo fator de autenticação}

Exemplo com conta Google + YubiKey:

\begin{enumerate}
  \item Compre uma YubiKey (recomenda-se comprar duas de uma vez: uma para uso diário, uma de backup, armazenadas em locais diferentes).
  \item Acesse myaccount.google.com Segurança Verificação em duas etapas.
  \item Selecione "Adicionar chave de segurança" insira a YubiKey na porta USB-C ou aproxime-a da área NFC do telefone.
  \item Toque o contato metálico da YubiKey conforme as instruções para concluir o registro.
  \item Recomenda-se registrar uma segunda chave de backup.
  \item Nas configurações de "Verificação em duas etapas", defina a chave de hardware como método preferencial e o código SMS como backup.
\end{enumerate}

Principais sites compatíveis: Google, GitHub, Microsoft, Apple ID, Facebook, Twitter, Dropbox, GitLab, AWS, Coinbase, Binance, entre outros.

\begin{tcbinfo}{Estratégia de duas chaves}
O maior risco da chave de hardware é “perdê-la”. Sugestões: (1) Registre duas chaves diferentes ao mesmo tempo; (2) Uma para carregar sempre (chaveiro), outra guardada em um cofre doméstico ou caixa de segurança bancária; (3) Alguns sites permitem imprimir o “código de recuperação de backup” para armazenamento offline. Assim, mesmo que a chave principal seja perdida, você não ficará permanentemente bloqueado fora da conta.

\end{tcbinfo}

\begin{tcbwarn}{Cuidado com o canal de compra}
Compre sempre na loja oficial ou em revendedores autorizados. Não compre chaves de segunda mão ou de origem desconhecida — em teoria, hardware pré-configurado pode ter backdoors implantados. O firmware da YubiKey não pode ser atualizado após a fabricação, mas a Yubico publica continuamente avisos de segurança; basta acompanhar o site oficial.

\end{tcbwarn}

\booksection{VIII. Compartilhamento de senhas em família/equipe}

No ambiente familiar, senhas de Netflix, Wi-Fi doméstico, painel do roteador e fechadura inteligente precisam ser compartilhadas por toda a família. No ambiente de equipe, senhas de servidores, bancos de dados e contas de redes sociais precisam de colaboração entre várias pessoas.\textbf{Enviar senhas pelo WeChat ou escrevê-las em documentos compartilhados são práticas de alto risco}— se o histórico de conversas ou a permissão do documento vazar, a senha também vaza.

\booksubsection{O “Cofre compartilhado” (Vault) do 1Password}

O 1Password é o produto de referência para compartilhamento familiar e em equipe. Seu mecanismo central são\textbf{múltiplos cofres}:

\begin{itemize}
  \item \textbf{Cofre pessoal}: Só você pode ver, como contas bancárias, e-mail e outras contas privadas.
  \item \textbf{Cofre compartilhado}: Familiares ou membros da equipe podem ver, como Netflix, Wi-Fi doméstico, cartão de crédito compartilhado.
  \item \textbf{Permissões graduadas}: A versão para equipes permite definir que certos membros tenham "apenas visualização" ou "apenas visualização sem edição".
\end{itemize}

Familiares ou membros da equipe não precisam compartilhar a senha mestra. Cada um faz login com sua própria conta, mas pode ver o conteúdo dos cofres compartilhados autorizados. Quando uma senha é alterada, todos os membros são sincronizados em tempo real.

\booksubsection{A função “Organização” do Bitwarden}

A versão gratuita do Bitwarden suporta compartilhamento de senhas pessoais (compartilhamento limitado por tempo de uma única senha via função “Send”). A versão paga permite criar “Organizações”, compartilhando itens com familiares ou colegas de equipe, com permissões detalhadas. O Bitwarden também oferece opção de auto-hospedagem, adequada para empresas com altas exigências de soberania de dados.

\booksubsection{Ambiente familiar: forma segura de compartilhar senhas de Netflix, Wi-Fi etc.}

\begin{itemize}
  \item \textbf{Netflix / sites de streaming}: Use o "Cofre compartilhado" do 1Password Família para armazenar a senha de login. Cada membro da família faz login individualmente, sem precisar enviar no grupo.
  \item \textbf{Senha do Wi-Fi doméstico}: Entre dispositivos iOS, é possível compartilhar a senha do Wi-Fi diretamente via AirDrop para visitantes, sem precisar ditá-la. Android 12+ também tem a função "Compartilhar Wi-Fi por proximidade".
  \item \textbf{Fechaduras inteligentes, eletrodomésticos inteligentes}: Sempre que possível, use a função "Membros da família" do aplicativo do fabricante para criar contas individuais para cada pessoa, em vez de compartilhar uma conta de administrador.
  \item \textbf{Painel do roteador doméstico}: Armazene a senha de administrador no cofre compartilhado do 1Password. Crie um SSID separado para o Wi-Fi de convidados.
\end{itemize}

\begin{tcbinfo}{Custo-benefício do 1Password Família}
O 1Password Families custa cerca de 4,99 USD/mês, suporta até 5 membros da família, cada um com seu próprio cofre privado + 4 cofres compartilhados. Isso dá menos de 1 USD/mês por pessoa, sendo atualmente a melhor solução de compartilhamento familiar. O Bitwarden Família custa 10 USD/ano para 6 pessoas, com custo-benefício ainda melhor, mas a experiência de interface é um pouco inferior.

\end{tcbinfo}

\booksubsection{Ambiente de equipe: compartilhamento de senhas de servidores/banco de dados}

\begin{itemize}
  \item \textbf{1Password Teams / Business}: Permite criar múltiplos cofres por departamento ou grupo de projeto, com controle granular de permissões e auditoria de todos os acessos.
  \item \textbf{Bitwarden Teams / Enterprise}: Solução de código aberto, com opção de auto-hospedagem. A versão Enterprise suporta integração SSO e sincronização de diretório.
  \item \textbf{Princípio de separação de responsabilidades}: A senha de administrador não deve ser conhecida por apenas uma pessoa — defina 2 a 3 "contatos de emergência" para garantir que o negócio não pare quando alguém essencial sair ou estiver de férias.
  \item \textbf{Processo de desligamento}: Ao sair de um funcionário, revogue imediatamente o acesso ao gerenciador de senhas e, se necessário, faça a rotação de todas as senhas sensíveis que ele tenha acessado.
\end{itemize}

\booksubsection{Acesso de Emergência (Emergency Access): se algo acontecer com você}

Imagine este cenário: você gerencia sozinho todas as contas importantes da família, mas sofre um acidente e é hospitalizado ou infelizmente falece. Sua família não consegue acessar contas bancárias, e-mails ou seguros.\textbf{A função “Acesso de Emergência”}foi projetada exatamente para isso — autorizar um familiar/amigo de confiança a obter acesso ao seu cofre de senhas quando você não puder responder.

\textbf{Configuração do Acesso de Emergência no Bitwarden:}

\begin{enumerate}
  \item Faça login no Bitwarden Web Vault Configurações Acesso de Emergência.
  \item Adicione o e-mail do contato de confiança e convide-o a aceitar.
  \item Defina o "período de espera" (recomendado 7-14 dias): a janela de tempo entre a solicitação da outra pessoa e sua confirmação.
  \item Se você não responder durante o período de espera (recusar ou não agir), ao expirar o prazo a pessoa recebe automaticamente o acesso.
  \item Você pode recusar manualmente a solicitação ou revogar a permissão de acesso de emergência a qualquer momento.
\end{enumerate}

\textbf{Acesso de Emergência no 1Password:}

\begin{enumerate}
  \item O 1Password não chama diretamente de "Acesso de Emergência", mas é possível implementar pela função "Organização Família/Equipe" — adicione a família à versão Família e coloque alguns itens no "Cofre compartilhado".
  \item Para usuários da versão pessoal, recomenda-se informar a senha mestra a um familiar de alta confiança ou escrevê-la em um envelope lacrado guardado em um cofre.
  \item O 1Password também fornece um "Código de recuperação". Ao ativar a autenticação de dois fatores, imprima e guarde-o.
\end{enumerate}

\begin{tcbpractice}{Configure agora: deixe uma "chave reserva" para sua família}
\begin{enumerate}
  \item Escolha um familiar de confiança como "contato de emergência".
  \item Configure o Acesso de Emergência no gerenciador de senhas (Bitwarden) ou adicione ao cofre compartilhado familiar (1Password).
  \item Anote a senha mestra + código de recuperação em papel, lacre em um envelope e guarde no cofre doméstico ou informe a um familiar confiável.
  \item Revise semestralmente para confirmar se o contato de emergência ainda é adequado.
\end{enumerate}

\end{tcbpractice}

\booksection{IX. Melhores práticas no uso do gerenciador de senhas}

\begin{itemize}
  \item \textbf{Defina uma senha mestra forte}: A senha mestra é a última linha de defesa. Defina uma senha longa e complexa e memorize-a.
  \item \textbf{Ative a autenticação de dois fatores}: Ative a autenticação de dois fatores para a conta do gerenciador de senhas, adicionando uma camada extra de proteção.
  \item \textbf{Verifique regularmente senhas fracas}: A maioria dos gerenciadores de senhas tem uma função de "auditoria de segurança" que ajuda a encontrar senhas fracas e duplicadas.
  \item \textbf{Não esqueça a senha mestra}: A senha mestra não pode ser recuperada. Esquecê-la significa perder todas as senhas. Recomenda-se anotá-la em papel e guardá-la em local seguro.
  \item \textbf{Acesso de Emergência}: Configure um contato de emergência. Se você não puder acessar, uma pessoa de confiança poderá obter suas senhas.
\end{itemize}

\begin{tcbwarn}{Aviso importante}
Embora o gerenciador de senhas seja conveniente, ele também tem riscos: se a senha mestra vazar, todas as senhas ficam expostas. Portanto, a senha mestra nunca deve ser igual a outras senhas, nem escrita em locais fáceis de encontrar. Além disso, recomenda-se escolher um provedor com bom histórico de segurança.

\end{tcbwarn}

\begin{tcbpractice}{Experimente agora: comece a usar o gerenciador de senhas}
\begin{enumerate}
  \item Escolha um gerenciador de senhas (recomenda-se Bitwarden ou o nativo do sistema).
  \item Registre uma conta e defina uma senha mestra forte.
  \item Importe as senhas existentes (se houver) ou comece a adicionar novas senhas.
  \item Instale a extensão do navegador ou o aplicativo móvel e experimente o preenchimento automático.
  \item Verifique o relatório de auditoria de segurança e encontre as senhas que precisam ser melhoradas.
\end{enumerate}

\end{tcbpractice}
